% GENERATED FILE. Edit canonical lessons, then run npm run generate:books.
% canonical-source-hash: fnv1a64:207937186d16b6e6
% canonical-lessons: RU-C17-no, RU-C17-a, RU-C17-odin, RU-C17-drugoy, RU-W08-tochka

\chapter{But, Twice Over}
\label{ch:ru-but-twice-over}
\hlchaptermodality{17}

\begin{chapteropening}
\textbf{By the end of this chapter:} \emph{I can contradict with no, set two facts side by side with a, count one, name the other one, divide a fact between two people with odin ... drugoy, and close a sentence with the three marks Russian writes only at the end.}

English hides two different jobs under one word, but. Russian splits them: no denies the first half, a sets two facts beside each other, and an English speaker reaches for the wrong one until the whereas test is stated out loud. drugoy then pays back a root the book has been holding since the friends chapter -- drug, the companion, is literally the other one. The chapter closes on the page itself: Russian, unlike Spanish, opens nothing and marks only the end.
\end{chapteropening}

\section[no]{\ru{но} — \textquotedblleft{}but,\textquotedblright{} said flatly}

\label{lesson:RU-C17-no}



\begin{quote}
One word from the drinks chapter, one from the verbs, and the two newest words in the book.

\begin{itemize}
\raggedright
  \item \emph{Recall:} \emph{\ru{любить}}, English \emph{love} itself, and the letter \textbf{\ru{л}} that appears from nowhere in \emph{\ru{я} \ru{люблю}}
  \item \emph{Recall:} \emph{\ru{кофе}}, which sailed Arabic to Turkish to Dutch to Russia, and the letter \textbf{\ru{ф}} it needed
  \item \emph{Recall:} \emph{\ru{и}}, and its old pointing stem, and \emph{\ru{тоже}} with \textbf{\ru{не}} after it
\end{itemize}
\end{quote}

\subsection*{You'll want to know first — \ru{но}}
\begin{itemize}
\raggedright
  \item \textbf{\ru{но}} — \emph{no} — \textbf{but}
\end{itemize}

Two letters, and it sits between the halves exactly where \textbf{\ru{и}} sat. Where \textbf{\ru{и}} adds, \textbf{\ru{но}} takes back:

\begin{itemize}
\raggedright
  \item \textbf{\ru{я} \ru{люблю} \ru{чай}, \ru{но} \ru{не} \ru{кофе}} — \emph{ya lyublyú chai, no ne kófe} — I like tea, \textbf{but not} coffee.
\end{itemize}

Notice the comma. Russian puts one in front of \textbf{\ru{но}} every time, and this is the first joining word in the book that has asked for one — \textbf{\ru{и}} did not.

\textbf{\ru{но}} contradicts. Whatever the first half led you to expect, the second half denies it, and nothing softer is meant.

\begin{cousinweb}
Almost every word this book has taken apart has handed you an English relative. \textbf{\ru{но}} does not. It goes back to a bare Slavic particle, written in the oldest manuscripts with a short vowel Russian has since dropped altogether, and the etymologists disagree about where even that came from.

That is worth saying out loud rather than papering over: a language's oldest, smallest, most-used words are often the ones whose history is least secure, because they are too short to leave much evidence behind.
\end{cousinweb}

\subsection*{Your turn}
\begin{itemize}
\raggedright
  \item \emph{Say it:} \emph{\ru{я} \ru{люблю} \ru{чай}, \ru{но} \ru{не} \ru{кофе}}
  \item \emph{Build:} the same sentence with \textbf{\ru{и}} instead of \textbf{\ru{но}}, and hear it stop contradicting
  \item \emph{Build:} two things you do, the second one denied, with \textbf{\ru{но} \ru{не}}
\end{itemize}

\begin{tcolorbox}[breakable,colback=teal!4,colframe=teal!35!black,title={Before you move on}]
Say \textquotedblleft{}but.\textquotedblright{} (That is \textbf{\ru{но}}.) Does a comma go in front of it? (\textbf{Yes} — always.) Did \textbf{\ru{и}} want one? (\textbf{No}.) Does \textbf{\ru{но}} have a secure English cousin? (\textbf{No} — the histories disagree.) And one look back: how do you say \textquotedblleft{}me neither\textquotedblright{}? (\textbf{\ru{я} \ru{тоже} \ru{не}} plus the verb.)
\end{tcolorbox}
\section[a]{\ru{а} — the other \textquotedblleft{}but,\textquotedblright{} and the one English does not have}

\label{lesson:RU-C17-a}



\begin{quote}
Two words from the food chapter, and the two joiners you already own.

\begin{itemize}
\raggedright
  \item \emph{Recall:} \emph{\ru{чай}}, which came overland from Mandarin through Persian while English \emph{tea} came by sea
  \item \emph{Recall:} \emph{\ru{хлеб}}, borrowed from Germanic before Russian was Russian, and the letter \textbf{\ru{х}} in it
  \item \emph{Recall:} \emph{\ru{и}} joining two whole clauses, and \emph{\ru{но}} taking the second one back
\end{itemize}
\end{quote}

\subsection*{You'll want to know first — \ru{а}}
\begin{itemize}
\raggedright
  \item \textbf{\ru{а}} — \emph{a} — \textbf{and, but} — the \emph{whereas} one
\end{itemize}

One letter again. Russian's two commonest joining words, \textbf{\ru{и}} and \textbf{\ru{а}}, are both single letters, and you learned to write both of them long before anyone said they were words.

\begin{itemize}
\raggedright
  \item \textbf{\ru{я} \ru{люблю} \ru{чай}, \ru{а} \ru{ты} \ru{любишь} \ru{кофе}} — \emph{ya lyublyú chai, a ty lyúbish kófe} — I like tea, \textbf{and you} like coffee.
\end{itemize}

Nothing is denied there. Two facts are set side by side so that the difference between them can be seen. English has no separate word for this and reaches for \emph{but}, or \emph{and}, or \emph{whereas}, depending on the sentence.

Like \textbf{\ru{но}}, \textbf{\ru{а}} takes a comma in front of it.

\begin{grammarlens}[title={Grammar lens — which of the two to reach for}]
Ask one question: \textbf{is the second half denying the first, or standing beside it?}

\begin{itemize}
\raggedright
  \item \textbf{\ru{но}} — denying. \emph{I like tea, \textbf{but not} coffee.}
  \item \textbf{\ru{а}} — standing beside. \emph{I like tea, \textbf{and you} like coffee.}
\end{itemize}

English hides both jobs under one word, so an English speaker will reach for \textbf{\ru{но}} far too often. The test that works: if you could say \emph{whereas}, use \textbf{\ru{а}}.
\end{grammarlens}

\subsection*{Your turn}
\begin{itemize}
\raggedright
  \item \emph{Say it:} \emph{\ru{я} \ru{люблю} \ru{чай}, \ru{а} \ru{ты} \ru{любишь} \ru{кофе}}
  \item \emph{Say it:} \emph{\ru{я} \ru{люблю} \ru{чай}, \ru{но} \ru{не} \ru{кофе}} — and hear the two sentences part company
  \item \emph{Build:} one of each, about you and somebody else
\end{itemize}

\begin{tcolorbox}[breakable,colback=teal!4,colframe=teal!35!black,title={Before you move on}]
Which word denies, and which sets two facts side by side? (\textbf{\ru{но}} denies; \textbf{\ru{а}} sets side by side.) Which English word can you swap in to test for \textbf{\ru{а}}? (\emph{Whereas}.) Which two joining words are single letters? (\textbf{\ru{и}} and \textbf{\ru{а}}.) And one look back: which drink came overland and which came by sea? (Tea came overland as \textbf{\ru{чай}}; English \emph{tea} came by sea.)
\end{tcolorbox}
\section[adín]{\ru{один} — \textquotedblleft{}one,\textquotedblright{} and five English words that are the same word}

\label{lesson:RU-C17-odin}



\begin{quote}
The pair of friends, and two joining words.

\begin{itemize}
\raggedright
  \item \emph{Recall:} \emph{\ru{друг}}, on a root meaning \textbf{to accompany}, whose English descendants have all died out
  \item \emph{Recall:} \emph{\ru{подруга}}, the same word with a prefix and the feminine \textbf{-\ru{а}}
  \item \emph{Recall:} \emph{\ru{или}}, built from \textbf{\ru{и}} plus \textbf{\ru{ли}}, and \emph{\ru{а}} the \emph{whereas} word
\end{itemize}
\end{quote}

\subsection*{You'll want to know first — \ru{один}}
\begin{itemize}
\raggedright
  \item \textbf{\ru{один}} — \emph{adín} — \textbf{one}
\end{itemize}

Stress the second syllable, and the unstressed \textbf{\ru{о}} at the front says \emph{a}, which is the reduction rule the letters chapter left you with: \emph{a-DÍN}, never \emph{o-DIN}.

\begin{itemize}
\raggedright
  \item \textbf{\ru{один} \ru{друг}} — \emph{adín drug} — \textbf{one friend}.
\end{itemize}

It is the first number this book has given you, and it is here for a reason: it is half of a joining pattern that arrives in the next lesson.

\begin{cousinweb}
Strip the strengthening piece off the front of \textbf{\ru{один}} and what is left goes back to PIE \textbf{*óynos}, \textquotedblleft{}one.\textquotedblright{} That single root gave English \textbf{five} separate everyday words:

\begin{itemize}
\raggedright
  \item \textbf{one} — the number itself
  \item \textbf{an} and \textbf{a} — which are \emph{one}, worn down to nothing
  \item \textbf{alone} — \emph{all-one}
  \item \textbf{only} — \emph{one-ly}
\end{itemize}

Latin took the same root and made \textbf{unus}, behind \emph{unit}, \emph{union} and \emph{universe}. Russian kept it in \textbf{\ru{один}} and, unlike English, never wore any of it down into an article — which is why Russian has no word for \emph{a} at all.
\end{cousinweb}

\subsection*{Your turn}
Say these aloud:

\begin{itemize}
\raggedright
  \item \emph{\ru{один}} — \emph{a-DÍN}, with the first vowel reduced
  \item \emph{\ru{один} \ru{друг}}
  \item the English family — \emph{one, an, a, alone, only} — all one root
\end{itemize}

\begin{tcolorbox}[breakable,colback=teal!4,colframe=teal!35!black,title={Before you move on}]
Say \textquotedblleft{}one.\textquotedblright{} (That is \textbf{\ru{один}}.) Which syllable is stressed, and what does the first vowel say? (\textbf{The second}; the first \textbf{\ru{о}} says \emph{a}.) Name three English words from the same root. (\emph{One}, \emph{an}, \emph{alone} — or \emph{a}, or \emph{only}.) And one look back: what does \emph{\ru{друг}}'s root mean? (\textbf{To accompany}.)
\end{tcolorbox}
\section[drugóy]{\ru{другой} — \textquotedblleft{}the other one,\textquotedblright{} and it is your word for friend}

\label{lesson:RU-C17-drugoy}



\begin{quote}
The pair of siblings, and yesterday's number.

\begin{itemize}
\raggedright
  \item \emph{Recall:} \emph{\ru{брат}}, one of the most secure cognate sets in the whole family — \emph{brother}, \emph{frater}, \emph{bhrātar}
  \item \emph{Recall:} \emph{\ru{сестра}}, which completes the pair just as securely
  \item \emph{Recall:} \emph{\ru{один}}, and \emph{\ru{ни} … \ru{ни} …} standing in front of both halves
\end{itemize}
\end{quote}

\subsection*{You'll want to know first — \ru{другой}}
\begin{itemize}
\raggedright
  \item \textbf{\ru{другой}} — \emph{drugóy} — \textbf{other, another}
\end{itemize}

Look at the front of it. It is \textbf{\ru{друг}}, your word for \emph{friend}, with an ending stitched on — and the meaning that came out is \emph{the other one}.

That is not a coincidence and it is not a pun. In the oldest Slavic, a \textbf{\ru{друг}} was a companion: the one who is not you, the second of a pair. Ask who \emph{the other one} is, and the answer is the person beside you.

\begin{itemize}
\raggedright
  \item \textbf{\ru{другой} \ru{брат}} — \emph{drugóy brat} — the \textbf{other} brother.
\end{itemize}

\begin{grammarlens}[title={Grammar lens — \ru{один} … \ru{другой} …}]
Put the two words at the front of two halves and you have the last of the joining patterns: the one that hands out a fact to each of two people.

\begin{itemize}
\raggedright
  \item \textbf{\ru{один} \ru{брат} \ru{любит} \ru{чай}, \ru{другой} \ru{любит} \ru{кофе}} — \emph{adín brat lyúbit chai, drugóy lyúbit kófe} — \textbf{one} brother likes tea, \textbf{the other} likes coffee.
\end{itemize}

Like \textbf{\ru{ни} … \ru{ни} …}, this pattern stands in \textbf{front} of its halves rather than between them. Unlike \textbf{\ru{ни} … \ru{ни} …}, it denies nothing: it divides.
\end{grammarlens}

\subsection*{Your turn}
\begin{itemize}
\raggedright
  \item \emph{Say it:} \emph{\ru{другой}} — and hear \emph{\ru{друг}} at the front of it
  \item \emph{Say it:} \emph{\ru{один} \ru{брат} \ru{любит} \ru{чай}, \ru{другой} \ru{любит} \ru{кофе}}
  \item \emph{Build:} the same split with two sisters, and two drinks of your own
\end{itemize}

\begin{tcolorbox}[breakable,colback=teal!4,colframe=teal!35!black,title={Before you move on}]
Which word for a person is hiding inside \textbf{\ru{другой}}? (It is \textbf{\ru{друг}}.) Why does \emph{companion} end up meaning \emph{the other one}? (\textbf{Because the companion is the one who is not you}.) Which two joining patterns stand in front of their halves rather than between them? (\textbf{\ru{ни} … \ru{ни} …} and \textbf{\ru{один} … \ru{другой} …}.)
\end{tcolorbox}
\section[tochka, vopros, vosklitsanie]{. ? ! — the three marks that close a Russian sentence}

\label{lesson:RU-W08-tochka}



\begin{quote}
Two nouns from earlier, before the page itself is the lesson.

\begin{itemize}
\raggedright
  \item \emph{Recall:} \emph{\ru{семья}}, a cousin of English \emph{home} through the root meaning \textbf{to settle}
  \item \emph{Recall:} \emph{\ru{ухо}}, inherited straight from the root behind English \emph{ear}
  \item \emph{Recall:} \emph{\ru{тоже}}, and how \textbf{\ru{не}} turns it into \emph{me neither}
  \item \emph{Recall:} four of the last eight letters — \textbf{\ru{ф}}, \textbf{\ru{ц}}, \textbf{\ru{й}}, \textbf{\ru{ю}} — and which of the eight was the false friend
\end{itemize}
\end{quote}

\subsection*{Writing — the three closing marks}
Russian uses the same three marks English does, in the same three shapes:

\begin{itemize}
\raggedright
  \item \textbf{.} — \textbf{\ru{точка}} — the full stop, and the Russian word means simply \emph{a dot}
  \item \textbf{?} — the question mark
  \item \textbf{!} — the exclamation mark
\end{itemize}

A statement ends with the dot. Nothing else is needed and nothing goes at the front:

\begin{quote}
\ru{брат} \ru{и} \ru{сестра}.
\end{quote}

Draw the dot on the writing line, not floating above it. That is the whole of the stroke, and it is the mark you will make more often than any other.

\begin{grammarlens}[title={Grammar lens — Russian closes, and never opens}]
Spanish opens a question with an upside-down mark and closes it with the ordinary one. Russian does \textbf{not}. One mark, at the end, and nothing at the beginning:

\begin{itemize}
\raggedright
  \item \textbf{\ru{чай} \ru{или} \ru{кофе}?} — one mark, at the end
  \item \textbf{\ru{ухо}!} — the same rule for the exclamation
\end{itemize}

So a Russian question looks exactly like a Russian statement until you reach the last character. What tells a listener it is a question is the \textbf{voice}, not the page — which is why the reading and the saying have to be practised together.
\end{grammarlens}

\subsection*{Your turn}
\begin{itemize}
\raggedright
  \item \emph{Write it:} \emph{\ru{брат} \ru{и} \ru{сестра}} and close it with a dot on the line
  \item \emph{Write it:} \emph{\ru{чай} \ru{или} \ru{кофе}} and close it with one question mark, nothing at the front
  \item \emph{Say it:} the same two words as a statement, then as a question, and hear the only difference
\end{itemize}

\begin{tcolorbox}[breakable,colback=teal!4,colframe=teal!35!black,title={Before you move on}]
How many marks does a Russian question take, and where? (\textbf{One}, at the end.) What does \textbf{\ru{точка}} mean as an ordinary word? (\textbf{A dot}.) What tells a listener a Russian sentence is a question? (\textbf{The voice}.) And one look back: which pattern hands one fact to each of two people? (\textbf{\ru{один} … \ru{другой} …})
\end{tcolorbox}
